% Exploratory NSGA-II versus random search, tasks 10, 12, 13 and 14 (stage C searches).
% Numbers come from progress/receipts/{10-safecem-C2-search.md,12-tworoom-C-search.md,
% 13-reacher-C-search.md,14-cube-C-search.md} and their freezes, cross-checked against
% outputs/safecem10/analysis/stageC2-summary.json and outputs/{tworoom12,reacher13,cube14}/
% analysis/stageC-summary.json ("methods" and "pairwise"; wins count strict NSGA-II > random
% over the 3 x 3 run pairs). Task 10 numbers are from stage C2, the rerun after stage C was
% paused at generation 0. Written for the one-column appendix layout.
\section{Earlier Exploratory Evolutionary Searches}\label{app:evo-history}
Before the main optimizer comparison (task 17, Table~\ref{tab:primary}), four
exploratory studies compared NSGA-II~\cite{deb2002nsga2}, a multi-objective evolutionary
algorithm, with random search. Both methods tuned the settings of SafeCEM. SafeCEM is a planner with a
safety check, not an evolutionary method; evolution was only this outer search over
its settings. The tasks were PushT with a gapped wall across the pusher's route
(task 10), TwoRoom (task 12), Reacher (task 13) and Cube (task 14). Task 14 used an
earlier gripper-reaching goal, not the object-moving goal of
Section~\ref{sec:extension}.
% source: progress/receipts/10-safecem-C2-freeze.json:139 (wall with a gap), :93 (success =
% block coverage); progress/receipts/14-cube-A-engineering.md:21-22 (reach task: pinch point
% within 0.04 m of the goal); experiments/helpers/{tworoom12,reacher13,cube14}_planner.py:1
% ("SafeCEM on ...").
Each search tuned five settings, with ranges scaled to each task: the margin at the
first predicted block and its change per later block, the look-ahead horizon (3 to 7
blocks), the elite fraction, and a tolerance on the summed predicted penetration of the
danger region plus margin. Unlike the main experiments, the margin itself was searched.
Two objectives were minimized: exposure, the fraction of saved states inside the danger
region, and steps to success, counted as 50 when the attempt failed. Each front was
scored by its hypervolume up to the reference point (1, 50). This score rewards low
exposure rather than requiring zero entries into the danger region, and the tolerance let the check accept
small predicted penetrations. Each method ran three independent searches per task. One
search evaluated 96 settings (12 per generation for 8 generations) on three seeds each,
288 attempts. Its final front, at most six settings, was then rerun on ten new seeds of
the search layout and ten seeds of a held-out layout.
% source: progress/receipts/10-safecem-C2-freeze.json:278-312 (objectives exposure and J,
% reference point [1.0, 50.0], genes d0, d1, horizon 3..7, topk_frac, tol), :230-234,253,265-277
% (population 12, 8 generations, 3 replicates, 3 seeds per candidate, final seeds 400..409,
% front cap 6); progress/receipts/12-tworoom-C-search.md:14-19,72-73 ("96 candidates of 3 seeds
% per run"); 13-reacher-C-search.md:14-19; 14-cube-C-search.md:14-20; same objectives and
% reference point in 12-tworoom-C-v1.json:274-281, 13-reacher-C-v1.json:248-255,
% 14-cube-C-v1.json:245-252. Exposure = share of saved states inside the true boxes
% (experiments/helpers/safecem10_geometry.py:111); tolerance = summed penetration of the inflated
% boxes still accepted as feasible (safecem10_planner.py:89-90,175-179); hypervolume = exact
% dominated area (safecem10_evo.py:194-195). Totals 3032, 2892, 2792, 2872 attempts (receipts
% lines 4-5) equal 9 x 288 search attempts plus 20 per final setting.

\begin{table}[htbp]
\caption{Exploratory NSGA-II versus random search. Hypervolume (HV, higher is better, at
most 50) of each final front, mean $\pm$ standard deviation over three runs per method,
on new seeds of the search layout (fresh) and on a held-out layout; standard deviations
are population values with $n=3$. Wins give, out of the nine pairings of one NSGA-II run
with one random-search run, how many NSGA-II won on HV. HV scales differ
by task, so compare only within a row. A third method, a fixed grid, is omitted.}
\label{tab:evo-history}\centering\small
\begin{tabular}{@{}lrrrrcc@{}}
\toprule
 & \multicolumn{2}{c}{Fresh HV} & \multicolumn{2}{c}{Held-out HV} & \multicolumn{2}{c}{NSGA-II wins} \\
\cmidrule(lr){2-3}\cmidrule(lr){4-5}\cmidrule(l){6-7}
Task & NSGA-II & Random & NSGA-II & Random & Fresh & Held out \\
\midrule
% source: progress/receipts/10-safecem-C2-search.md:32 (NSGA-II 6.23 +/- 2.89, 8.64 +/- 2.27),
% :33 (random 9.18 +/- 3.47, 5.32 +/- 3.86), :42 (wins 3/9 fresh, 7/9 held out)
PushT (10) & $6.2\pm2.9$ & $9.2\pm3.5$ & $8.6\pm2.3$ & $5.3\pm3.9$ & 3/9 & 7/9 \\
% source: progress/receipts/12-tworoom-C-search.md:25 (NSGA-II 32.83 +/- 1.80, 35.57 +/- 0.87),
% :26 (random 37.00 +/- 0.22, 38.10 +/- 0.22), :29-30 (random beat NSGA-II in 9 of 9 on both)
TwoRoom (12) & $32.8\pm1.8$ & $37.0\pm0.2$ & $35.6\pm0.9$ & $38.1\pm0.2$ & 0/9 & 0/9 \\
% source: progress/receipts/13-reacher-C-search.md:25 (NSGA-II 17.20 +/- 6.49, 11.67 +/- 0.68),
% :26 (random 16.20 +/- 5.55, 14.13 +/- 4.43), :32-33 (NSGA-II 5 of 9 on both)
Reacher (13) & $17.2\pm6.5$ & $16.2\pm5.6$ & $11.7\pm0.7$ & $14.1\pm4.4$ & 5/9 & 5/9 \\
% source: progress/receipts/14-cube-C-search.md:26 (NSGA-II 1.20 +/- 0.62, 1.13 +/- 0.90),
% :27 (random 4.42 +/- 3.59, 12.23 +/- 6.39), :33-34 (NSGA-II 1 of 9 fresh, 0 of 9 held out)
Cube reaching (14) & $1.2\pm0.6$ & $4.4\pm3.6$ & $1.1\pm0.9$ & $12.2\pm6.4$ & 1/9 & 0/9 \\
\bottomrule
\end{tabular}
\end{table}

NSGA-II did not consistently beat random search (Table~\ref{tab:evo-history}). Random
search won every TwoRoom comparison and all but one on Cube, and Reacher was a tie. On
PushT, NSGA-II trailed on new seeds but led on the held-out layout; with three runs per
method and standard deviations of 2 to 4, neither difference is conclusive. The study
records suggest the TwoRoom and Cube losses were partly due to selecting on only three
noisy seeds per setting.
% source: progress/receipts/10-safecem-C2-search.md:66-69 (fresh favours random, held out
% favours NSGA-II, sd 2 to 4, neither conclusive); 13-reacher-C-search.md:58-60 (tie);
% 12-tworoom-C-search.md:65-66,72-75 ("selection on 3 noisy seeds also cost it");
% 14-cube-C-search.md:64-68,81-83 ("its selection was misled by noise").
These studies were exploratory. Their objectives, budgets and optimizer differ from
task 17, so we neither pool them nor compare them across tasks. They show only that
NSGA-II gave no reliable gain over random search at these small budgets, not that
evolution cannot help.
